\begin{claim}{D.43}{Proved}{Global orientation of punctured-torus involutions}
Every involution in the full once-punctured-torus loop-outer
group has the global orientation multiplier $1$ or $-1$,
with the same sign as its $2$-adic determinant.
Consequently the orientation character of every finite
subgroup generated by involutions has values in the global
two-element subgroup $\{1,-1\}\subset\widehat\Z^\times$.
\end{claim}
\begin{proof}
Write $F=\widehat F(x,y)$ and $c=[x,y]$ for the torus
cusp.  D.27 supplies an actual order-two representative
$\alpha$ of the given outer involution.  Its peripheral
multiplier is
\[
 \lambda=\det(\alpha|_{F^{\rm ab}})\in\widehat\Z^\times,
                    \qquad \lambda^2=1,
                                                        \tag{D.43a}
\]
by the exact Heisenberg comparison of D.33.
Thus every $\lambda_p$ is $1$ or $-1$.

\emph{The characteristic Heisenberg cover.}
For an odd prime $q$, take the characteristic quotient
\[
 F\longrightarrow G_q,\qquad
 G_q=\langle X,Y,Z\mid X^q=Y^q=Z^q=1,\
              [X,Y]=Z,\ [X,Z]=[Y,Z]=1\rangle,
                                                        \tag{D.43b}
\]
with $x\mapsto X$, $y\mapsto Y$, $c\mapsto Z$.
Collection of these words identifies $G_q$ with the
order-$q^3$ Heisenberg group.  Explicitly it is
$\F_q^2\times\F_q$, with multiplication
\[
 (v,z)(w,t)=(v+w,z+t+\tfrac12\omega(v,w)),
 \qquad \omega((a,b),(a',b'))=ab'-ba'.
                                                        \tag{D.43c}
\]
Take $X=((1,0),0)$, $Y=((0,1),0)$ and $Z=((0,0),1)$.
The kernel $N_q$ is the closed subgroup generated by
the laws $\gamma_3F$ and all $q$th powers, so it is
characteristic.  More precisely it is their closed
normal closure: the exponent-$q$, class-two quotient
has the normal forms $X^aY^bZ^d$ in (D.43b), and
(D.43c) verifies that all $q^3$ forms are distinct.
This proves the characteristic-quotient identification.

The corresponding discrete cover of a once-punctured
torus has $q^2$ cusps, one for each coset of
$\langle Z\rangle$ in $G_q$, each of degree $q$.
Fill all these cusps.  Completing the finite-cover
presentation gives the exact quotient
\[
 N_q\longrightarrow
 \Pi_q=N_q/
   \overline{\langle\!\langle h c^q h^{-1}:
                       h\in F\rangle\!\rangle}_{N_q}
       \cong\widehat\pi_{g_q,0},
 \qquad 2g_q=2+q^3-q^2 .
                                                        \tag{D.43d}
\]
Representatives for the $q^2$ cosets suffice in the
normal closure.  The genus formula follows by multiplying
the punctured torus's Euler characteristic by $q^3$
and then adding its $q^2$ filled discs.
The original deck group induces its actual homology
representation on this closed surface.
Since $\alpha$ preserves $N_q$ and permutes all the
filled peripheral classes, it induces an actual
involution $\beta_q$ of $\Pi_q$.

For every prime $p$, put $V_p=H_1(\Pi_q,\Q_p)$,
let $\rho$ be its $G_q$ representation, and let
$A$ be the map induced by $\beta_q$.  If $\theta$
is the automorphism induced by $\alpha$ on $G_q$,
the precise deck comparison is
\[
 A^2=1,\qquad
 A\rho(h)A^{-1}=\rho(\theta(h)),\qquad
                    \theta(Z)=Z^{\lambda\bmod q}.
                                                        \tag{D.43e}
\]
Indeed representatives in $F$ for the two sides differ
only by an element of $N_q$, whose conjugation is inner
on $\Pi_q$ and hence trivial on its first homology.

The nondegenerate alternating surface pairing satisfies
\[
 \langle Av,Aw\rangle=\lambda_p\langle v,w\rangle,
       \qquad
 \langle\rho(h)v,\rho(h)w\rangle=\langle v,w\rangle .
                                                        \tag{D.43f}
\]
Here is an exact verification of the first multiplier.
The involution permutes the odd number $q^2$ of lifted
cusps, so one cusp class is fixed.  Fill every other
cusp first.  This is an invariant closed normal
quotient, and gives a surface group of genus $g_q$
with one cusp, which $\alpha$ preserves with power
$\lambda$.  Filling its last cusp does not change
first homology.  The degree-two boundary bivector
calculation in (D.41h) therefore proves the first
identity in (D.43f) for precisely $A$ on $V_p$.
The second identity comes from the orientation-preserving
geometric deck maps of the specified finite surface cover.
Thus no realization of $\alpha$ as a homeomorphism
has been used.

We compute the whole deck module by finite cover chains.
The cover before filling retracts onto the regular
cover of a two-loop rose, so its first homology has
representation class $1+\Q_p[G_q]$.
The relative second homology of the attached discs
is $\Q_p[G_q/\langle Z\rangle]$.
The exact homology sequence of this filling, including
the invariant orientation class in degree two, gives
\[
 [V_p]=
       2[1]+[\Q_p[G_q]]
                 -[\Q_p[G_q/\langle Z\rangle]].
                                                        \tag{D.43g}
\]
All these maps are deck-equivariant.  They compute
the homology of the completed surface groups through
the finite-index and surface-homology comparisons
already used in A.2 and D.42.

Over a characteristic-zero field containing a primitive
$q$th root $\zeta$, let $V_j$ be the subspace where $Z$
acts as $\zeta^j$, for $j\in\F_q^\times$.
The last permutation module in (D.43g) is trivial on
$Z$, and the regular representation has each of its
$q$ central eigenvalues with multiplicity $q^2$.
Semisimplicity therefore gives
\[
                         \dim V_j=q^2 .
                                                        \tag{D.43h}
\]
The $G_q$-invariant alternating form pairs $V_j$
only with $V_{-j}$, and gives a nondegenerate
pairing between them.

\emph{One negative odd-prime sign determines every sign.}
Suppose $\lambda_q=-1$ at an odd prime $q$.
For any prime $p$, extend scalars to $\Q_p(\zeta)$.
By (D.43e), $A$ exchanges $V_j$ and $V_{-j}$.
Choose bases of these two spaces and write $P$ for
their nonsingular pairing matrix, $B$ for
$A:V_j\to V_{-j}$ and $C$ for $A:V_{-j}\to V_j$.
Then $CB=BC=1$, and alternation in (D.43f) gives
\[
                  -B^{\mathsf t}P^{\mathsf t}C
                                      =\lambda_p P.
                                                        \tag{D.43i}
\]
The dimension $q^2$ is odd.  Taking determinants
in (D.43i) yields $-1=\lambda_p^{q^2}$, hence
$\lambda_p=-1$ by (D.43a).
This works also at $p=q$ and $p=2$, since the
calculation is over characteristic-zero fields.
Thus
\[
       \lambda_q=-1\text{ for an odd }q
                  \quad\Longrightarrow\quad\lambda=-1.
                                                        \tag{D.43j}
\]

\emph{The positive $3$-adic case.}
It remains to prove that $\lambda_3=1$ forces
$\lambda_2=1$.  The matrix of $\alpha$ on
$H_1(F,\Z_3)$ has square one and determinant one.
Its two eigenvalues over $\Q_3$ are consequently
both $1$ or both $-1$, and it is diagonalizable,
so the matrix is $I$ or $-I$.

In the group $G_3$ of (D.43c), the resulting $\theta$
has the form
\[
                \theta(v,z)=(Mv,z+\ell(v)),
                      \qquad M=I\text{ or }-I,
                                                        \tag{D.43k}
\]
with $\ell$ a linear functional.
For $M=I$, $\theta^2=1$ implies $2\ell=0$ and
$\ell=0$.
For $M=-I$, conjugation of $\theta$ by the inner
automorphism of $(u,0)$ replaces $\ell(v)$ by
$\ell(v)-2\omega(u,v)$.  Choose $u$ with
$\omega(u,v)=\ell(v)/2$.  Lifting $(u,0)$ to $F$
and conjugating the actual automorphism $\alpha$
by that inner automorphism preserves its order two
and its peripheral multiplier.  It gives exactly
\[
             \theta=1\quad\hbox{or}\quad
             \theta=\iota,\qquad\iota(v,z)=(-v,z).
                                                        \tag{D.43l}
\]
This is conjugation of an actual automorphism,
not multiplication by an arbitrary inner lift.

Work now over the unramified quadratic field
$L=\Q_2(\zeta_3)$, with
$\sigma(\zeta_3)=\zeta_3^{-1}$.
For $j=1,2$ the irreducible $G_3$ module $W_j$
with central character $\zeta_3^j$ has basis
$e_r$, $r\in\F_3$, and matrices
\[
 X e_r=e_{r+1},\qquad
 Y e_r=\zeta_3^{-jr}e_r,\qquad
 Z e_r=\zeta_3^j e_r .
                                                        \tag{D.43m}
\]
These matrices satisfy (D.43b).
The distinct eigenvalues of $Y$ and the cyclic
action of $X$ show irreducibility.  The diagonal
projectors for $Y$, followed by powers of $X$,
span the full matrix algebra.  Hence the central
block of $L[G_3]$ is $M_3(L)$.
Together with (D.43g)--(D.43h), this gives
\[
                       V_j=W_j\otimes_L U_j,
                             \qquad\dim U_j=3.
                                                        \tag{D.43n}
\]

The matrix $J=I$ implements $\theta=1$ on $W_j$.
For $\theta=\iota$, take $J e_r=e_{-r}$.
This rational permutation matrix implements
$X\mapsto X^{-1}$, $Y\mapsto Y^{-1}$, fixes $Z$,
and has $J^2=I$ and $J^{\mathsf t}J=I$.
Thus (D.43e) and the full matrix-algebra commutant
give
\[
                A|_{V_j}=J\otimes B_j,\qquad B_j^2=I.
                                                        \tag{D.43o}
\]
Choose a multiplicity basis for $U_1$ and its
$\sigma$-image for $U_2$.
Since the original representation and $A$ are
defined over $\Q_2$, we then have
$B_2=\sigma(B_1)$.
The canonical invariant pairing between $W_1$
and $W_2$ is the ordinary dot product in (D.43m).
Consequently the nondegenerate cross-pairing in
(D.43f) has the form $I_3\otimes S$ for an
invertible $3$-by-$3$ matrix $S$ on the multiplicity
spaces.  Equation (D.43f) becomes
\[
                 B_1^{\mathsf t}S\,\sigma(B_1)
                                      =\lambda_2 S.
                                                        \tag{D.43p}
\]
Because $B_1^2=I$, its determinant is $1$ or $-1$,
both fixed by $\sigma$.
Taking determinants in (D.43p) therefore gives
$1=\lambda_2^3$.  Together with $\lambda_2^2=1$,
this proves $\lambda_2=1$.

If $\lambda_2=-1$, the last conclusion forces
$\lambda_3=-1$, so (D.43j) gives $\lambda=-1$.
If $\lambda_2=1$, (D.43j) excludes every negative
odd-prime coordinate, and $\lambda=1$.
This proves the claim.  Multiplicativity proves
the stated consequence for subgroups generated
by involutions; no conclusion about their
conjugacy is inferred.

The program \texttt{code/part\_c/torus\_involution\_d43.py}
checks the characteristic Heisenberg presentations,
the finite-cover chain ranks and central-character
dimensions, and every $G_3$ involution with determinant
one in (D.43k), including the exact inner-conjugacy
formula in (D.43l).  It also verifies the
Schr\"odinger matrices and their rational implementer.
The field and determinant arguments above prove
the assertion for the full profinite groups.
\end{proof}
